\section{The uniform Fourier enclosure}\label{v9:fourier:section}
Fix an actual index $n\ge N\ge42$, a centered variance-one marginal $X$
and raw threshold $H$. Their values may vary with $n$: the estimates below are applied anew to
that law and threshold, using common ceilings. In particular, no limiting
law or comparison of characteristic functions at different indices is used.
Write
\[
 f(t)=\E e^{itX},\quad\gamma=\E X^3,\quad
 e_n(t)=e^{-nt^2/2}(1-in\gamma t^3/6),\quad V_n=f^n-e_n,
\]
\[
 I_0=\int_0^{1/2}\frac{|\Im(e^{-itH}V_n)|}{t^3}dt,
 \qquad I_\infty=\int_{1/2}^\infty\frac{|V_n|}{t^3}dt.
\]

The phase $e^{-itH}$ retains the raw threshold. At low frequency we use
the imaginary part, whose cancellation is stronger than an absolute-value
bound on $V_n$; on the tail an absolute-value bound is sufficient.

\paragraph{Compensated inversion.}
\begin{lemma}\label{v9:fourier:inversion}
For a finite signed measure $\mu$ of mass and first moment zero with
$\int x^2d|\mu|<\infty$, every $H\in\R$ satisfies
\[
 \int(x-H)|x-H|d\mu(x)
 =-\frac4\pi\int_0^\infty\frac{\Im(e^{-itH}\widehat\mu(t))}{t^3}dt,
\]
with absolute convergence.
\end{lemma}
\begin{proof}
First take $y>0$ and substitute $v=ty$. Integration by parts gives
\[
 \int_0^\infty\frac{ty-\sin ty}{t^3}dt
 =\frac{y^2}{2}\int_0^\infty\frac{1-\cos v}{v^2}dv
 =\frac\pi4y^2.
\]
The boundary term vanishes: $v-\sin v=O(v^3)$ at zero and $O(v)$ at
infinity. To evaluate the constant, insert
$v^{-2}=\int_0^\infty s e^{-sv}ds$. Positive Tonelli gives
\[
 \int_0^\infty\frac{1-\cos v}{v^2}dv
 =\int_0^\infty s\{s^{-1}-s/(s^2+1)\}ds=\pi/2.
\]
Oddness extends the first identity to $\pi y|y|/4$ for every real $y$.
Moreover $ty-\sin ty$ has the sign of $y$. Thus, for $y=x-H$, the
absolute double integral against $d|\mu|(x)dt$ equals
$\pi\int(x-H)^2d|\mu|/4<\infty$.
Absolute Fubini is now justified. Since
$\int(x-H)d\mu=0$, the linear compensation disappears and gives the
stated Fourier integral. This argument uses measures directly: atoms,
including an atom at $H$, require no continuity convention.
\end{proof}
Apply this to the sum law minus the scale-$\sqrt n$ pushforward of
$\phi(x)\{1+\gamma(x^3-3x)/(6\sqrt n)\}dx$.
The comparator has mass one, mean zero and transform $e_n$.
For $z=H/\sqrt n$, the identity $q_H(\sqrt n x)=nq_z(x)$ fixes the
normalization. Three Gaussian integrations by parts, interpreted
distributionally using $q_z'''=4\delta_z$, give
\[
 \int q_z(x)(x^3-3x)\phi(x)dx=4\phi(z).
\]
Consequently the signed measure in the inversion lemma integrates $q_H$
to $n\varepsilon$, with exactly the $\varepsilon$ of the core identity.
Multiplication by $-\sqrt n/2$ therefore gives its remainder
$E_n=-n^{3/2}\varepsilon/2$ and the bound:
\begin{equation}\label{v9:fourier:ledger}
 E_n=\frac{2\sqrt n}{\pi}\int_0^\infty
       \frac{\Im(e^{-itH}V_n)}{t^3}dt,
 \qquad |E_n|\le\frac{2\sqrt n}{\pi}(I_0+I_\infty).
\end{equation}

\paragraph{The complete low-frequency bound.}
Assume common actual ceilings
$\mu_4\le M\le111/10$, $|\gamma|\le g\le37/20$, $\mu_6\le M_6$,
$|(3-\mu_4)/24|\le A$, $|(10\gamma-\E X^5)/120|\le Z$.
Put
\begin{gather*}
 E=M_6/720+g^2/72+(1/2+g^2/144)^3/5,\qquad B=A+Z/2+E/4,\\
 a_0=(N-1)/(2N),\quad b_0=(45-M)(1-2/N)/96,\quad
 \Lambda_N=\sqrt{\pi/2}(N/(N-1))^{3/2}.
\end{gather*}
\begin{proposition}\label{v9:fourier:low}
For $c_0\ge\Lambda_N$, $c_1\ge\Lambda_N/2$,
\[
 I_0\le\frac{L_0+L_H|H|}{\sqrt n}+\frac{L_2}{n},
\]
where
\[
 \begin{gathered}
 L_0=c_0Z+c_1Ag+251g^3/17280,\quad L_H=c_0A+251g^2/4800,\\
 L_2=E/(2a_0^2)+B^2/(2b_0^3).
 \end{gathered}
\]
\end{proposition}
\begin{proof}
For $h=e^{-t^2/2-i\gamma t^3/6}$ and $0\le t\le1/2$, Taylor's integral
remainder gives
\begin{equation}\label{v9:fourier:taylor}
 f-h=-\frac{3-\mu_4}{24}t^4-i\frac{10\gamma-\E X^5}{120}t^5+r_6,
 \quad |r_6|\le Et^6,\quad |f-h|\le Bt^4.
\end{equation}
For $w=-t^2/2-i\gamma t^3/6$, use
$|w|\le t^2(1/2+g^2/144)<1/6$ and
$|e^w-1-w-w^2/2|\le|w|^3/5$; the sixth-degree term of $w^2/2$
costs $g^2t^6/72$, and the marginal remainder costs $\mu_6t^6/720$.
Four integrations of $1-\cos u\ge0$ give
$\cos u\le1-u^2/2+u^4/24$. Symmetrization implies
\[
 |f|^2\le1-t^2+(M+3)t^4/12\le1-(45-M)t^2/48.
\]
Both $f,h$, and their segment, thus have modulus at most
$e^{-(45-M)t^2/96}$. Taylor expansion of the polynomial $v^n$ gives
\[
 f^n-h^n=n(f-h)h^{n-1}+R_2,\quad
 |R_2|\le\tfrac12n(n-1)B^2t^8e^{-(n-2)(45-M)t^2/96}.
\]
Integrating $r_6$ and $R_2$ after division by $t^3$ costs $L_2/n$:
$n-1\ge2a_0n$, $(n-2)(45-M)/96\ge b_0n$, and
$\int_0^\infty t^3e^{-At^2}dt=1/(2A^2)$,
$\int_0^\infty t^5e^{-At^2}dt=A^{-3}$.
It remains to integrate the fourth- and fifth-order terms in
\eqref{v9:fourier:taylor}; taking absolute values too early would lose
the fourth-order cancellation. Put $b=n-1$. After rotation the real
fourth-order coefficient contributes a sine of $Ht+b\gamma t^3/6$,
whereas the fifth-order coefficient contributes a cosine. Hence their
imaginary part is at most
\[
 ne^{-bt^2/2}\{(Z+A|H|)t^5+bAg t^7/6\}.
\]
After division by $t^3$, use
\[
 \int_0^\infty t^2e^{-bt^2/2}dt=\sqrt{\pi/2}\,b^{-3/2},\qquad
 \int_0^\infty t^4e^{-bt^2/2}dt=3\sqrt{\pi/2}\,b^{-5/2}.
\]
These produce $\Lambda_n(Z+A|H|+Ag/2)/\sqrt n$.
Since $n/(n-1)\le N/(N-1)$, this is bounded by the stated expression
with $\Lambda_N$.
Finally, putting $u=n\gamma t^3/6$ gives
\[
 |\Im(e^{-itH}(h^n-e_n))|
 \le e^{-nt^2/2}(|u|^3/6+|H|tu^2/2).
\]
Gaussian moments yield
$\sqrt{\pi/2}\{5g^3/432+g^2|H|/24\}/\sqrt n$;
$\sqrt{\pi/2}<251/200$ proves the stated coefficients.
All enlarged integrands are nonnegative.
\end{proof}

The sharper fifth-order input follows from an exact moment identity.
For acquired support $|X|\le R\le6$, $R\ge1$, and $\E|X|^5\le E_5$,
take $M_6=RE_5$ and $A=\max(2,M-3)/24$. For $Q=X^2-\gamma X-1$,
\[
 \E Q^2=\mu_4-1-\gamma^2,\qquad
 \E X^5=\gamma^3+2\gamma+\E[Q^2(X+2\gamma)].
\]
Reflect this auxiliary moment calculation so $t=|\gamma|$.
For $r\ge\max(R,2g)$ both signs of $\E X^5-10\gamma$ are bounded by
\[
 P_r=r(M-1)+(10-2M)t-rt^2+t^3,\quad
 Q_r=r(M-1)+(2M-10)t-rt^2-t^3.
\]
The coefficients of $\mu_4$ are $r\mp2t\ge0$.
If $M\le5$, $P_r\ge Q_r$ and
$P_r\le r(M-1)+(10-2M)t-(r-g)t^2$.
If $M>5$, $P_r\le r(M-1)$ and
$Q_r\le r(M-1)+(2M-10)t-rt^2$.
The quadratic maxima use vertices clipped to $[0,g]$.
For $r=6$, $P_6''=-12+6t<0$: when $M\le5$ and $P_6'(g)\ge0$
use $P_6(g)$; otherwise use $6(M-1)+(10-2M)^2/[4(6-g)]$.
For $M>5$ also use $6(M-1)+(2M-10)^2/24$.
Minimizing the applicable ceilings with $E_5+10g$ and dividing by $120$
therefore gives the claimed valid $Z$.

\paragraph{Every high-frequency piece and its floor gate.}
For two iid copies assume an actual event $\mathcal E$ with
$\Pp(\mathcal E)\ge\alpha>0$, $D=|X-X'|\ge d>0$ on $\mathcal E$,
and $\E[D^2\mid\mathcal E]\le M_E$.
Fix $b=9/10$, $0<r\le1/4$, and set
\begin{gather*}
 h=333/106-(22/7)^3r^2/6>0,\quad\kappa=\alpha(23/48)^2d^2,
 \quad L(x)=(x+2)/(x^2+4x+2),\\
 S_3=\sum_{j=1}^{20}j^{-3}+1/800,\quad
 S_5=\sum_{j=1}^{20}(j-r)^{-5}+[4(20-r)^4]^{-1},\\
 c(u)=u/2-(M+3)u^2/24,\quad c_* =\min\{c(1/4),c(81/100)\},\\
 c_q=1/8-(M+3)/384,\qquad s=[1-(M+3)53/600]/2.
\end{gather*}
\begin{proposition}\label{v9:fourier:high}
If $c_*>0$, $\sqrt n I_\infty\le\sum_{i=1}^6T_i(n)$, where
\begin{align*}
 T_1&=2\sqrt n e^{-nc_*},&
 T_2&=\frac{M_E}{4\pi^{3/2}h\sqrt\alpha}
          \{S_3+3S_5/(n\alpha h^2)\},\\
 T_3&=\frac{M_E\sqrt n}{2\pi^2}e^{-n\alpha h^2r^2},&
 T_4&=\frac{e^{-n\kappa b^2}}{2b^4\kappa\sqrt n},\\
 T_5&=\frac{\sqrt n}{2b^2}e^{-n\alpha(23/48)^2},&
 T_6&=\sqrt n e^{-n/8}\{2L(n/8)+\sqrt{M-1}/3\}.
\end{align*}
The optional replacements, minimized with the respective simple terms, are
\[
 \widetilde T_1=2\sqrt n e^{-nc_q}L(ns/4)\quad(s>0),\qquad
 \widetilde T_4=\frac{\sqrt n e^{-n\kappa b^2}L(n\kappa b^2)}{2b^2}.
\]
The common-floor sum at $N$ is valid if
\[
 2Nc'>1\quad(c'\in\{c_*,\alpha h^2r^2,\alpha(23/48)^2,1/8\}).
\]
The optional first term requires $2Nc_q>1$, and the optional fourth
$N\kappa b^2\ge1/4$. Without an optional gate, retain the simple term.
\end{proposition}
\begin{proof}
We first bound the characteristic function on $[1/2,b]$, then partition
the remaining half-line around its possible lattice peaks. The comparator
is integrated over the entire half-line separately. The six terms record
these disjoint purposes: the initial interval, nonzero-cell cores and
wings, two parts of the zero cell, and the comparator.

On $[1/2,b]$,
$|f|^n\le e^{-nc(t^2)}\le e^{-nc_*}$ and $\int t^{-3}dt<2$,
giving $T_1$. Its optional chord follows from
$c(u)-c_q-s(u-1/4)=(M+3)(u-1/4)(81/100-u)/24\ge0$.
For $x=nab^2>0$ the integration formula is
\begin{equation}\label{v9:fourier:tail}
 \sqrt n\int_b^\infty e^{-na t^2}t^{-3}dt
 \le\frac{\sqrt n e^{-x}}{2b^2}L(x).
\end{equation}
Substitution and positive Fubini give the left side as this prefactor
times $1-xQ(x)$, where $Q(x)=\int_0^\infty e^{-u}/(x+u)du$.
For $P(u)=u^2-4u+2$ and $D_x=P(-x)=x^2+4x+2$, use
$P(u)-D_x=(u+x)(u-x-4)$ and
$\int_0^\infty P(u)e^{-u}du=\int_0^\infty uP(u)e^{-u}du=0$.
The two zero integrals follow from $\int_0^\infty u^k e^{-u}du=k!$.
For clarity, the two divisions are
\[
 \int_0^\infty\frac{P(u)e^{-u}}{x+u}du=D_xQ(x)-x-3,
\]
\[
 \int_0^\infty\frac{P(u)^2e^{-u}}{x+u}du
 =D_x\int_0^\infty\frac{P(u)e^{-u}}{x+u}du.
\]
Indeed, the remainder in the second equality is the integral of
$P(u)(u-x-4)$, which vanishes by those two zero moments. Thus
\[
 Q(x)-\frac{x+3}{D_x}
 =\frac1{D_x^2}\int_0^\infty\frac{P(u)^2e^{-u}}{x+u}du\ge0.
\]
This positive remainder proves $1-xQ(x)\le L(x)$ for every $x>0$,
not just for large $x$. It proves \eqref{v9:fourier:tail}; applying that
formula to the chord majorant, extended to $[1/2,\infty)$, proves the
optional first term.

For $t\ge b$, symmetrization and convex Jensen give
\[
 |f(t)|^n\le e^{-n\E\sin^2(t(X-X')/2)}
 \le\E[e^{-n\alpha\sin^2(tD/2)}\mid\mathcal E].
\]
Here $|f|^2=1-2\E\sin^2(t(X-X')/2)$, so $1-v\le e^{-v}$ first
gives the exponential. Restricting the expectation to $\mathcal E$ and
using its probability at least $\alpha$ then gives the conditional Jensen
bound. Positive Tonelli permits a separate partition for each $D$:
with $\ell=2\pi/D$, use $[0,\ell/2]$ and nonzero cells
$t=\ell(j+v)$, $j\ge1$, $|v|\le1/2$.
This distance-dependent partition is made inside the conditional integral;
it does not assume a common lattice for the original law.
On $|v|\le r$, $|\sin\pi v|\ge h|v|$; on the wings it is at least $hr$,
by $\sin x\ge x-x^3/6$ and monotonicity on $[0,\pi/2]$.
Taylor's remainder gives
\[
 \{(j-v)^{-3}+(j+v)^{-3}\}/2\le j^{-3}+6(j-r)^{-5}v^2.
\]
Pair $v$ with $-v$ and extend each nonnegative Gaussian core integral to
$\R$. With $a=n\alpha h^2$, its two required moments are
$\int_{\R}e^{-av^2}dv=\sqrt{\pi/a}$ and
$\int_{\R}v^2e^{-av^2}dv=\sqrt\pi/(2a^{3/2})$.
The Jacobian is $t^{-3}dt=D^2(j+v)^{-3}dv/(4\pi^2)$. After multiplying
by $\sqrt n$ and averaging $D^2$, these factors give
\[
 \frac{M_E}{4\pi^{3/2}h\sqrt\alpha}
 \sum_{j\ge1}\{j^{-3}+3(j-r)^{-5}/(n\alpha h^2)\}.
\]
For decreasing summands the tails beyond $20$ are bounded by their
integrals from $20$ to infinity. These are $1/800$ and
$[4(20-r)^4]^{-1}$, respectively, proving $T_2$. The wing weight is at most
$\int_{\ell/2}^\infty t^{-3}dt=D^2/(2\pi^2)$, giving $T_3$.
Extending these sets below $b$ only increases their bounds.
In the zero cell, $u=tD/2\le\pi/2$. On $u\le1/2$,
$\sin u\ge23u/24$ and $D\ge d$. Integrate over the enlarged
$[b,\infty)$ with $t^{-3}\le t/b^4$ to obtain $T_4$;
\eqref{v9:fourier:tail} gives its optional form.
On $u\ge1/2$, $\sin u\ge23/48$ gives $T_5$.
Cell endpoints have zero Lebesgue measure even when $D$ has atoms.
For the comparator,
$|e_n|\le e^{-nt^2/2}(1+n|\gamma|t^3/6)$.
Use \eqref{v9:fourier:tail} at $b=1/2$ and
$\int_{1/2}^\infty e^{-nt^2/2}dt\le2e^{-n/8}/n$;
$\gamma^2\le M-1$ from $\E(X^2-\gamma X-1)^2\ge0$ gives $T_6$.
No comparator or lattice tail is omitted.

For floor evaluation, $\sqrt n e^{-nc'}$ decreases when $2nc'>1$,
whereas $n^{-1}$ and $n^{-1/2}e^{-nc'}$ decrease for $c'>0$.
Also $L'(x)=-(x^2+4x+6)/(x^2+4x+2)^2<0$.
The derivative of $\sqrt xe^{-x}L(x)$ has the sign of
$4-10x-22x^2-13x^3-2x^4$, negative and decreasing for $x\ge1/4$.
These prove every stated gate on a common scalar majorant, not by
comparing characteristic functions of different laws.
\end{proof}

For any complete HIGH sum $K_H$ at $N$, the final bound is
\begin{equation}\label{v9:fourier:conclusion}
 |E_n|\le\frac2\pi\{L_0+K_H+L_2/\sqrt N+L_H|H|\}
 \le\frac{16}{25}\{L_0+K_H+L_2/\sqrt N+L_H|H|\}.
\end{equation}
Complete HIGH windows may be minimized; the LOW triple stays complete.
Round denominator roots down and numerator roots up, using the constants
below. For $x\ge0$, choose rational $0\le y\le\min(x,64)$ and round
$1/\sum_{j=0}^{64}y^j/j!$ upward to majorize $e^{-x}$.
The floor proof concerns the analytical functions, not monotonicity of
this rounded exponential helper.